\subsection{Unitary characters of a locally compact abelian group}

DEFINITION 1. --- A unitary character of G is a continuous homomorphism from G into the multiplicative group U of complex numbers of modulus 1.

In other words, a unitary character is a continuous function \(\chi\) on G, with complex values, such that:

\[
\chi (x y) = \chi (x) \chi (y), \qquad | \chi (x) | = 1 \qquad (x, y \in \mathrm{G}).
\]

In this chapter, we will often simply say “character” instead of “unitary character.”

Let E be a Hilbert space of dimension 1, and let \(\chi\) be a unitary character of G. The map that assigns to \(x \in \mathrm{G}\) the homothety of ratio \(\chi(x)\) in E is a continuous linear isometric representation of G on E. Conversely, every bounded continuous linear representation of G on E is obtained by this process, and in particular is unitary.

It is immediate that the product of two unitary characters, the inverse of a unitary character, and the constant function equal to 1 are unitary characters. Consequently, the set \(\widehat{\mathbf{G}}\) of unitary characters of G forms a group under multiplication. This group is commutative. On the other hand, the map \((\chi_1,\chi_2)\mapsto \chi_1\chi_2^{-1} = \chi_1\overline{\chi}_2\) is continuous for the compact-open topology, and \(\widehat{\mathbf{G}}\), endowed with the topology of compact convergence, is a topological group (TG, X, p.~6, corollary 2 and remark 1).

DEFINITION 2. --- The topological group \(\widehat{\mathbf{G}}\) is called the dual group of G.

Since G is locally compact, the map \((x,\chi)\mapsto\chi(x)\) is continuous on \(G\times\widehat{G}\) (TG, X, p.~28, th. 3).

Recall that \(\mathcal{M}^{1}(G)\) denotes the unital involutive Banach algebra of bounded complex measures on G (Example 4 of I, p.~99). For every bounded complex measure \(\mu\in\mathcal{M}^{1}(G)\) and every \(\chi\in\widehat{G}\) , we denote by

\[
\chi (\mu) = \int_ {\mathrm{G}} \chi (x) d \mu (x)\tag{1}
\]

(cf. INT, VIII, §2, no. 6).

Lemma 1. --- For every \(\chi \in \widehat{\mathbf{G}}\), the map \(\mu \mapsto \chi(\mu)\) is a Hermitian character of the involutive Banach algebra \(\mathcal{M}^1(\mathbf{G})\).

By INT, VIII, §3, no. 3, prop. 11, the map \(\mu \mapsto \chi(\mu)\) is a character of the involutive Banach algebra \(\mathcal{M}^{1}(G)\). Moreover, we have:

\[
\chi (\mu^ {*}) = \int_ {\mathrm{G}} \chi (x ^ {- 1}) d \overline {{\mu}} (x) = \int_ {\mathrm{G}} \overline {{\chi (x)}} d \overline {{\mu}} (x) = \overline {{\chi (\mu)}}
\]

and this character is therefore Hermitian.

Thus we have defined a mapping from \(\widehat{G}\) into \(\mathsf{X}(\mathcal{M}^{1}(\mathrm{G}))\); it will be called canonical.

Let \(\chi \in \widehat{\mathbf{G}}\). The restriction of \(\mu \mapsto \chi(\mu)\) to \(\mathrm{L}^1(\mathrm{G})\) is nonzero (cf. INT, VIII, §2, no. 7, prop. 10). By restriction to the involutive Banach subalgebra \(\mathrm{L}^1(\mathrm{G})\), we therefore obtain a map from \(\widehat{\mathbf{G}}\) into \(\mathsf{X}(\mathrm{L}^1(\mathrm{G}))\), said to be canonical. It associates to \(\chi \in \widehat{\mathbf{G}}\) the Hermitian character

\[
f \mapsto \chi (f) = \chi (f \cdot d x) = \int_ {\mathrm{G}} f (x) \chi (x) d x \qquad (f \in \mathrm{L} ^ {1} (\mathrm{G}))\tag{2}
\]

of \(L^{1}(G)\) .

PROPOSITION 1. --- The canonical map from \(\widehat{\mathbf{G}}\) into \(\mathsf{X}(\mathrm{L}^{1}(\mathbf{G}))\) is a homeomorphism.

Let us denote this canonical map by ev and, for \(\chi \in \widehat{\mathbf{G}}\), let \(\mathrm{ev}_{\chi}\) be the image of the character \(\chi\) by ev, that is, the Hermitian character \(f \mapsto \chi(f)\) of \(\mathrm{L}^1(\mathrm{G})\). Considered as a mapping from \(\widehat{\mathbf{G}}\) into the dual of \(\mathrm{L}^1(\mathrm{G})\), the map ev is the composition of the injection of \(\widehat{\mathbf{G}}\) into \(\mathrm{L}^{\infty}(\mathrm{G})\), equipped with the weak topology \(\sigma(\mathrm{L}^{\infty}(\mathrm{G}), \mathrm{L}^{1}(\mathrm{G}))\), and the injection of \(\mathrm{L}^{\infty}(\mathrm{G})\) into the dual of \(\mathrm{L}^1(\mathrm{G})\), equipped with the topology of pointwise convergence. Since \(\widehat{\mathbf{G}}\) is a bounded subset of \(\mathrm{L}^{\infty}(\mathrm{G})\), the first map is continuous by Lebesgue's theorem (INT, IV, §3, no. 7, th. 6). The second is also continuous, by definition. This proves that the map ev is continuous.

If \(\chi \in \widehat{\mathrm{G}}\) and if \(f \in \mathrm{L}^1(\mathrm{G})\) , one has \(\mathrm{ev}_{\chi}(\varepsilon_x * f) = \chi(x)\mathrm{ev}_{\chi}(f)\) . Taking \(f\) such that \(\mathrm{ev}_{\chi}(f) \neq 0\) , we deduce that the map ev is injective.

Let \(\zeta\in\mathsf{X}(\mathrm{L}^{1}(\mathrm{G}))\) and let \(f\in\mathrm{L}^{1}(\mathrm{G})\) such that \(\zeta(f)\neq0\) . Let us define a map \(\chi\colon\mathrm{G}\to\mathbf{C}\) by setting, for \(x\in\mathrm{G}\) :

\[
\chi (x) = \frac {\zeta (\varepsilon_ {x} * f)}{\zeta (f)}.\tag{3}
\]

We have \(\chi(e) = 1\). Since the map \(x \mapsto \varepsilon_x * f = \gamma(x)(f)\) from G to L \(^1\) (G) is continuous (INT, VIII, §2, n° 5, prop. 8), the map \(\chi\) is continuous. It is bounded because, for all \(x \in \mathbf{G}\), one has

\[
| \chi (x) | \leqslant \frac {\| \varepsilon_ {x} * f \|}{| \zeta (f) |} = \frac {\| f \|}{| \zeta (f) |}
\]

(th. 1 of I, p.~29 and INT, VIII, loc. cit.).

Now let \(\mathfrak{B}\) be a base of the filter of neighborhoods of \(e\) consisting of compact neighborhoods. For every \(V \in \mathfrak{B}\), let \(g_V\) be a positive continuous function, vanishing outside \(V\) and with integral equal to 1 (lemma 1 of II, p.~200). For every function \(h \in L^1(G)\), we then have

\[
\varepsilon_ {x} * h = \lim _ {\mathrm{V}, \mathfrak {B}} (\varepsilon_ {x} * h) * g _ {\mathrm{V}} = \lim _ {\mathrm{V}, \mathfrak {B}} (\varepsilon_ {x} * g _ {\mathrm{V}} * h),
\]

in L \(^{1}\) (G), the limit being taken along the filter of sections of \(\mathfrak{B}\) (INT, VIII, §4, n° 7, prop. 20). In particular, since \(\zeta(\varepsilon_x * g_V * f) = \zeta(\varepsilon_x * g_V)\zeta(f)\), one deduces that

\[
\chi (x) = \lim _ {\mathrm{V}, \mathfrak {B}} \zeta (\varepsilon_ {x} * g _ {\mathrm{V}}).
\]

For every \(h\in \mathrm{L}^1 (\mathrm{G})\) , we obtain

\[
\zeta (\varepsilon_ {x} * h) = \lim _ {\mathrm{V}, \mathfrak {B}} \zeta (\varepsilon_ {x} * g _ {\mathrm{V}} * h) = \zeta (h) \lim _ {\mathrm{V}, \mathfrak {B}} \zeta (\varepsilon_ {x} * g _ {\mathrm{V}}) = \chi (x) \zeta (h).
\]

Consequently, we have, for \(x, y \in G\):

\[
\chi (x y) = \frac {\zeta (\varepsilon_ {x} * \varepsilon_ {y} * f)}{\zeta (f)} = \frac {\chi (x) \zeta (\varepsilon_ {y} * f)}{\zeta (f)} = \chi (x) \chi (y),
\]

which proves that \(\chi\) is a homomorphism from G to \(\mathbf{C}^*\). Since \(\chi\) is bounded and continuous, it is a unitary character of G. Moreover, if \(g \in \mathrm{L}^{1}(\mathrm{G})\), one has

\[
g * f = \int_ {\mathrm{G}} (\varepsilon_ {x} * f) g (x) d x
\]

in L \(^{1}\) (G) (INT, VIII, §1, n° 5, prop. 7), whence

\[
\begin{array}{c} \zeta (g) \zeta (f) = \zeta (g * f) = \int_ {\mathrm{G}} \zeta (\varepsilon_ {x} * f) g (x) d x \\ = \zeta (f) \int_ {\mathrm{G}} \chi (x) g (x) d x = \mathrm{ev} _ {\chi} (g) \zeta (f) \end{array}
\]

(INT, VI, §1, n°1, prop. 1), which shows that \(\zeta = \mathrm{ev}_{\chi}\) . Consequently, ev is surjective, hence bijective.

Finally let us show that the inverse map \(ev^{-1}\) is continuous. Let \(\zeta \in \mathsf{X}(\mathrm{L}^{1}(\mathrm{G}))\). Let \(f \in \mathrm{L}^{1}(\mathrm{G})\) be a function such that \(\zeta(f) \neq 0\). The set W of \(\xi \in \mathsf{X}(\mathrm{L}^{1}(\mathrm{G}))\) such that \(\xi(f) \neq 0\) is an open neighborhood of \(\zeta\) in \(\mathsf{X}(\mathrm{L}^{1}(\mathrm{G}))\). For every \(\xi \in W\), what precedes shows that \(\mathrm{ev}^{-1}(\xi)\) is the character

\[
x \mapsto \frac {\xi (\varepsilon_ {x} * f)}{\xi (f)}.
\]

Let \(\mathfrak{F}\) be a filter on \(W \subset X(L^{1}(G))\) converging to \(\zeta\). Since the set \(X(L^{1}(G))\) is bounded, hence equicontinuous, in \(L^{\infty}(G)\), the uniform structure of simple convergence coincides with the uniform structure of compact convergence (TG, X, p.~16, th. 1). Let K be a compact subset of G. The set of \(\varepsilon_{x} * f\) for \(x \in K\) is compact in \(L^{1}(G)\) (INT, VIII, §2, n° 5, prop. 8). We therefore have

\[
\lim _ {\xi , \mathfrak {F}} \xi (\varepsilon_ {x} * f) = \zeta (\varepsilon_ {x} * f)
\]

uniformly for \(x \in K\) . From this we deduce that

\[
\lim _ {\xi , \mathfrak {F}} \mathrm{ev} ^ {- 1} (\xi) = \mathrm{ev} ^ {- 1} (\zeta),
\]

hence ev \(^{-1}\) is continuous at \(\zeta\) . This completes the proof of the proposition.

From now on we identify a unitary character \(\chi\) of G with the character \(f\mapsto\int_{\mathrm{G}}f(x)\chi(x)dx\) of \(\mathrm{L}^{1}(\mathrm{G})\).

Remarks. --- 1) *The bijectivity of the map from \(\widehat{G}\) onto \(X(L^{1}(G))\) in prop. 1 is a special case of the correspondence between continuous representations of a locally compact group H (not necessarily commutative) and continuous representations of the algebra \(L^{1}(H)\).*

\begin{enumerate}
\def\labelenumi{\arabic{enumi})}
\setcounter{enumi}{1}
\tightlist
\item
  The canonical map from \(\widehat{G}\) into \(\mathsf{X}(\mathcal{M}^{1}(\mathsf{G}))\) is not surjective in general (II, p.~308, exerc. 14).
\end{enumerate}

COROLLARY 1. --- Every character of L \(^{1}\) (G) is Hermitian. The canonical map from X(Stell(G)) (def. 9 of I, p.~125) into X(L \(^{1}\) (G)) is a homeomorphism.

The first assertion follows from Proposition 1 and Lemma 1, restricted to \(L^{1}(G)\). The second follows from the first and from the corollary of Prop. 20 of I, p.~124.

Corollary 2. --- The topological group \(\widehat{\mathbf{G}}\) is locally compact.

Indeed, \(\mathsf{X}(\mathrm{L}^{1}(\mathrm{G}))\) is locally compact (corollary of Theorem 1 of I, p.~29).

We shall identify \(\widehat{G}\) with \(\mathsf{X}(\mathrm{L}^1 (\mathrm{G}))\) and \(\mathsf{X}(\mathrm{Stell}(\mathrm{G}))\). For \(x\in G\) and \(\chi \in \widehat{G}\), we shall denote by \(\langle \chi ,x\rangle\) the complex number \(\chi (x)\), which belongs to \(\mathbf{U}\).

We say that \(x\) and \(\chi\) are orthogonal if \(\langle \chi, x \rangle = 1\). Let A be a subset of G (resp. of \(\widehat{\mathbf{G}}\)); the set of elements of \(\widehat{\mathbf{G}}\) (resp. of G) orthogonal to A is a closed subgroup of \(\widehat{\mathbf{G}}\) (resp. of G) which is called the orthogonal of A and is denoted \(\mathrm{A}^{\perp}\). The orthogonal of G consists only of \(e\).

For \(x \in G\), let \(\eta(x)\) be the map from \(\widehat{G}\) to U defined by \(\chi \mapsto \langle \chi, x \rangle\). By definition of multiplication in \(\widehat{G}\), the map \(\eta(x)\) is a group homomorphism. It is continuous since the map \((x, \chi) \mapsto \langle \chi, x \rangle\) from \(G \times \widehat{G}\) to U is continuous (TG, X, p.~28, th. 3). We have thus defined a mapping \(\eta\), called canonical, from G into the bidual group \(\widehat{\widehat{G}}\); it is a group homomorphism. Moreover, the map \(\eta\) is continuous (TG, X, p.~28, th. 3). We shall prove later (II, p.~220, th. 2) that \(\eta\) is a topological group isomorphism from G onto \(\widehat{\widehat{G}}\).

Let G and H be locally compact abelian groups, and \(\varphi\colon\mathrm{G}\to\mathrm{H}\) a morphism of topological groups. For every \(\chi\in\widehat{\mathrm{H}}\), the map \(\chi\circ\varphi\) is a character of G, denoted \(\widehat{\varphi}(\chi)\). This definition translates to the formula

\[
\langle \chi , \varphi (x) \rangle = \langle \widehat {\varphi} (\chi), x \rangle\tag{4}
\]

for all \(\chi\in\widehat{H}\) and \(x\in G\). From this we deduce that \(\widehat{\varphi}\) is a morphism from the topological group \(\widehat{H}\) to the topological group \(\widehat{G}\); we say that \(\widehat{\varphi}\) is the dual of the morphism \(\varphi\).

Let K be a locally compact commutative group and \(\psi\colon\mathrm{H}\to\mathrm{K}\) a morphism of topological groups. The definition shows that \(\widehat{\psi\circ\varphi}=\widehat{\varphi}\circ\widehat{\psi}\). If \(\varphi\) is the identity map of G, then \(\widehat{\varphi}\) is the identity map of \(\widehat{\mathrm{G}}\). In particular, if \(\varphi\) is an isomorphism of topological groups, then so is \(\widehat{\varphi}\), and \(\widehat{\varphi}^{-1}\) is the dual of \(\varphi^{-1}\).

Lemma 2. --- Let G and H be locally compact abelian groups, and let \(f\colon\mathrm{H}\to\mathrm{G}\) be a morphism of topological groups. The kernel of \(\widehat{f}\) is the orthogonal of the image of \(f\).

By definition, we have \(\chi \in \mathrm{Ker}(\widehat{f})\) if and only if the restriction of \(\chi\) to the image of \(f\) is trivial.

PROPOSITION 2. --- Let \(n \geqslant 0\) be an integer and let \(\mathrm{G}_{1}, \ldots, \mathrm{G}_{n}\) be locally compact commutative groups. Let G be the product group of the groups \(\mathrm{G}_{j}\) for \(1 \leqslant j \leqslant n\). For \(1 \leqslant j \leqslant n\), let \(\lambda_{j}\) be the injection of \(\mathrm{G}_{j}\) into G which associates to \(x \in \mathrm{G}_{j}\) the element \((x_{k})\) such that \(x_{k} = e\) if \(k \neq j\) and \(x_{j} = x\). The map

\[
(\widehat {\lambda} _ {j}) _ {1 \leqslant j \leqslant n} \colon \widehat {\mathrm{G}} \to \prod_ {1 \leqslant j \leqslant n} \widehat {\mathrm{G}} _ {j}
\]

is an isomorphism of topological groups.

Let \(m\) be the product map \(\widehat{\mathbf{G}}^n\) to \(\widehat{\mathbf{G}}\), and, for every \(j\) such that \(1 \leqslant j \leqslant n\), let \(\pi_j\) be the projection of \(\mathbf{G}\) onto \(\mathbf{G}_j\). Let \(\mu\) be the morphism of topological groups \(m \circ (\widehat{\pi}_j)_j\) from \(\prod \widehat{\mathbf{G}}_j\) to \(\widehat{\mathbf{G}}\), so that

\[
\langle \mu ((\chi_ {j})), (x _ {j}) \rangle = \prod_ {j = 1} ^ {n} \langle \chi_ {j}, x _ {j} \rangle .
\]

The map \(\mu\) is continuous, and one verifies that \(\mu\) and \((\widehat{\lambda}_j)\) are inverse bijections of each other. The proposition follows.

Remark. --- The computation of the dual group of an infinite product of compact commutative groups is the subject of Corollary 4 of II, p.~234 below. The case of a locally compact commutative group that is an arbitrary product of locally compact groups follows from these two statements, since in such a product all but finitely many factors are compact (TG, I, p.~66, Prop. 14, b)).

